\subsection{The fixed norm module and full arithmetic sections}

\begin{claim}{AC.45}{Proved: section origin is weaker than point origin}
\label{clm:ACnormsectionorigin}
Let $C/k$ be smooth, proper and geometrically connected of genus
at least two over a finite field of characteristic $p$, and let
$P\in C(k)$. Use the full etale fundamental groups in the
first row below, and put
\[
 \Omega=\pi_1(C_{\bar k},\bar P),\qquad
 \Delta=\Omega^{(p')},\qquad
 N=\ker(\Omega\longrightarrow\Delta),\qquad
 V=N^{\mathrm{ab},p}.
                                                        \tag{AC.45.1}
\]
Here $N^{\mathrm{ab},p}$ is the maximal abelian pro-$p$
quotient. The actual point section $s_P$ acts on these
groups; write $F$ for arithmetic Frobenius and also for
its induced automorphism of $V$. For the pointed norm
module $M_P$ of AR.11 there is a natural exact sequence
of compact abelian pro-$p$ groups
\[
 0\longrightarrow V\xrightarrow{\ F-1\ }V
      \xrightarrow{\ \varepsilon\ }M_P^F
      \longrightarrow0.
                                                        \tag{AC.45.2}
\]

Every $m\in M_P^F$ is the compatible $p$-adic
abelianized difference from $s_P$ of a continuous section
of the \emph{full} arithmetic etale fundamental group
of $C$. That section projects to exactly $s_P$ in the
geometrically prime-to-$p$ quotient. If $m\ne0$, A1
implies that this full section is not geometrically
conjugate to the section of any point of $C(k)$.

Projection is surjective at every $F$-stable pointed
level, in particular
\[
                    M_P^F\twoheadrightarrow J_C(k)[p^\infty].
                                                        \tag{AC.45.8}
\]
Thus nonzero section-origin classes actually exist
whenever $J_C(k)[p^\infty]\ne0$.

The statement supplies no element with free deck orbit
at every level. In particular it neither proves nor
refutes A10(c) or the nontrivial-cover existence
assertion A3.
\end{claim}

\begin{proof}
\emph{1. The integral norm limit.}
Take the open normal subgroups $H\triangleleft\Delta$
that are stable under $F$. They are cofinal: every
finite geometric cover has a model over a finite
constant extension, hence its defining subgroup has
a finite Frobenius orbit, whose intersection is a
stable normal refinement. Let $U_H\subset\Omega$
be the inverse image of $H$. The point section
gives a model $D_H/k$ with compatible rational lift
$P_H$. Write
\[
                    T_H=T_pJ_{D_H}.
\]
The Abel map based at $P_H$ identifies
$U_H^{\mathrm{ab},p}$ with $T_H$. This comparison
includes the characteristic prime: see
\cite[Proposition~9.1, Lemmas~9.4--9.5 and
Corollary~9.6]{AC_MilneJac1986}.
If $H'\subset H$, the map induced by
$U_{H'}\hookrightarrow U_H$ is the norm map on
Tate modules. Indeed the pointed Abel maps commute
with the morphism of curves and its norm homomorphism
on Jacobians. Equivalently, its dual is pullback on
first etale cohomology, not the transfer multiplied
by the covering degree.

There is an $F$-equivariant isomorphism
\[
                 V\xrightarrow{\sim}\varprojlim_H T_H.
                                                        \tag{AC.45.3}
\]
Here is an elementary continuity verification. A
continuous character $c:N\to\mathbf Q_p/\mathbf Z_p$
has finite image. Choose an open normal subgroup
$O\triangleleft\Omega$ with $N\cap O\subset\ker c$.
Since $[O,N]\subset O\cap N$, the formula
\[
               \widetilde c(no)=c(n),\qquad n\in N, o\in O,
\]
defines a character on the open subgroup $NO$.
Its quotient $\Omega/(NO)$ has prime-to-$p$ order,
since it is a finite quotient of $\Delta$.
Restricting to a sufficiently small $F$-stable normal
subgroup containing $N$ therefore extends $c$ to
one of the $U_H$. Conversely, a character on $U_H$
vanishing on $N$ factors through the pro-prime-to-$p$
group $H$ and is zero. Thus
\[
 \operatorname{Hom}_{\mathrm{cont}}(N,\mathbf Q_p/\mathbf Z_p)
  =\varinjlim_H
    \operatorname{Hom}_{\mathrm{cont}}(U_H,\mathbf Q_p/\mathbf Z_p).
\]
Pontryagin duality gives (AC.45.3), with the transition
maps already identified as norms. No comparison with
the full prime-to-$p$ configuration fiber is involved.

\emph{2. Frobenius coinvariants and rational torsion.}
For every $H$ there is an exact sequence
\[
 0\longrightarrow T_H\xrightarrow{\ F-1\ }T_H
     \xrightarrow{\ \varepsilon_H\ }
       J_{D_H}(k)[p^\infty]\longrightarrow0.
                                                        \tag{AC.45.4}
\]
To check it, put $W_H=T_H\otimes_{\mathbf Z_p}\mathbf Q_p$.
The quotient $W_H/T_H$ is the group of geometric
etale $p$-primary points of $J_{D_H}$. The fixed
subgroup of $T_H$ is zero: otherwise it would give
a nonzero compatible infinitely $p$-divisible
system in the finite group $J_{D_H}(k)$.
Since $T_H$ has finite rank, $F-1$ is therefore
invertible on $W_H$, and
\[
 \varepsilon_H(a)=(F-1)^{-1}a\pmod {T_H}
                                                        \tag{AC.45.5}
\]
gives (AC.45.4). Its image is precisely the
$F$-fixed subgroup of $W_H/T_H$, which is
$J_{D_H}(k)[p^\infty]$. This is the full $p$-adic
Kummer comparison over a finite field; it does not
treat multiplication by $p$ as an etale isogeny.
The maps in (AC.45.4) commute with norms.

For the cofinal stable levels, being fixed by the
action on the entire pointed inverse system means
exactly that each coordinate is rational over $k$.
Consequently
\[
           M_P^F=\varprojlim_H J_{D_H}(k)[p^\infty].
                                                        \tag{AC.45.6}
\]
This also equips $M_P^F$ with its compact pro-$p$
topology, regardless of the unbounded-exponent issue
for the complete module $M_P$.

We verify the inverse-limit exactness explicitly.
For a compatible family $m_H$ on the right of
(AC.45.4), each fiber
$\varepsilon_H^{-1}(m_H)\subset T_H$ is nonempty
and compact. Any finite collection of norm
compatibility equations can be satisfied by choosing
a preimage at one common refinement and taking its
norms to the finitely many levels in the equations.
The remaining coordinates may be chosen arbitrarily
in their nonempty fibers. The finite-intersection
property gives a compatible lift in
$\varprojlim_H T_H$. This proves surjectivity in
(AC.45.2), without a uniform denominator or an
unjustified assertion about restricted transition maps.

If a compatible family $a_H$ maps to zero, write
$a_H=(F-1)b_H$ using (AC.45.4). Each $b_H$ is
unique because $F-1$ is injective on $T_H$.
Commutation with norms then makes the $b_H$
compatible. This proves the kernel equality in
(AC.45.2). Its left injectivity follows at each
level as well. Using (AC.45.3) completes the proof
of that exact sequence.

The asserted surjectivity of projection has an
additional elementary proof. For a transition
$D_{H'}\to D_H$ over $k$, norm composed with
pullback is multiplication by its degree $d$.
Since $d$ is prime to $p$, multiplication by $d$
is an automorphism of $J_{D_H}(k)[p^\infty]$.
Norm on these rational $p$-primary groups is
therefore split surjective. Their finite inverse
system has surjective projection to every
coordinate, proving (AC.45.8). No assertion that
freeness survives these projections is made.

\emph{3. Realization by a full arithmetic section.}
Let $m\in M_P^F$. Choose $v\in V$ with
$\varepsilon(v)=m$, and then choose $n\in N$
mapping to $v$. In the full arithmetic group
\[
        1\longrightarrow\Omega\longrightarrow\pi_1(C,\bar P)
          \longrightarrow G_k\longrightarrow1,
\]
let $F=s_P(\operatorname{Frob}_k)$ and form
$nF$. Every element of a profinite group defines
a continuous procyclic power homomorphism from
$\widehat{\mathbf Z}$. The power homomorphism for
$nF$, followed by the projection to
$G_k\simeq\widehat{\mathbf Z}$, is the identity:
it sends the dense generator $1$ to $1$.
It therefore defines a continuous section $s_n$.
Its prime-to-$p$ geometric projection is exactly
$s_P$, since $n\in N$.

At the pointed level $D_H$, the element $n$ belongs
to $U_H$. With the additive convention for
$T_H$, the abelianized cocycle difference satisfies
\[
      (s_n-s_P)(\operatorname{Frob}_k)=a_H,
                                                        \tag{AC.45.7}
\]
where $a_H$ is the image of $v$ in $T_H$.
Our Kummer convention is
$\delta(x)(F)=F(\widetilde x)-\widetilde x$;
thus its cohomology class is
$a_H\bmod(F-1)T_H$, corresponding to $m_H$
under (AC.45.5). This fixes the sign and verifies
the asserted compatible section realization.

\emph{4. Point origin and the centralizer.}
Suppose $m\ne0$ and $s_n$ were geometrically
conjugate to the section of $Q\in C(k)$.
Projection to the geometrically prime-to-$p$
group and A1 imply $Q=P$.
If $\omega\in\Omega$ conjugates $s_n$ to $s_P$,
its image in $\Delta$ centralizes $s_P(G_k)$,
because these two sections have identical
prime-to-$p$ projections. That geometric
centralizer is trivial by
\cite[Lemma~1.7]{AR_ST2009}. Hence $\omega\in N$.
Its conjugation changes the abelianized difference
by an element of $(F-1)V$, so (AC.45.2) would
give $m=0$, a contradiction.

Thus nonzero norm-module classes have actual full
section realizations that are not point sections.
The argument says nothing further about their
deck stabilizers. In particular it does not replace
the missing point-origin assertion by a freeness
assertion or manufacture a counterexample to A3.
\end{proof}

\paragraph{Relation to the cuspidal proof.}
The point-recognition argument in AG.40 starts from
finite neighborhoods that are known to contain
actual rational points. A section $s_n$ constructed
above does not carry that property merely because
it is a continuous splitting. Consequently its
existence supplies no input to that recognition
step. The exact sequence distinguishes the
Frobenius-fixed Tate module, which is zero, from
the Frobenius coinvariants, which form the
rational torsion norm module. Confusing those two
groups would incorrectly remove the A10 gap.

\paragraph{A small actual-curve example.}
Over $\mathbf F_2$, take the smooth proper model of
\[
                   C:\quad y^2+y=x+\frac1x+\frac1{x+1}.
                                                        \tag{AC.45.9}
\]
The three simple poles $0,1,\infty$ give genus two
and two-rank two by the Artin--Schreier genus and
Deuring--Shafarevich formulas. Each pole supplies
one rational ramification point. There are no
other $\mathbf F_2$-points. Over $\mathbf F_4$,
the two remaining possible $x$-values give
Artin--Schreier right sides of absolute trace one,
so there are still exactly three points. Hence
\[
 \#C(\mathbf F_2)=\#C(\mathbf F_4)=3,
 \qquad P_C(T)=1-T^2+4T^4,
 \qquad \#J_C(\mathbf F_2)=4.
                                                        \tag{AC.45.10}
\]
The exact elementary program
\texttt{norm\_module\_small\_curve.py} verifies
both point counts and the polynomial. In fact
$J_C(\mathbf F_2)\simeq(\mathbf Z/2)^2$:
writing the three ramification points as
$W_0,W_1,W_\infty$, the divisors of $x$ and
$x+1$ are $2(W_0-W_\infty)$ and
$2(W_1-W_\infty)$. The corresponding classes
are distinct and nonzero, because a difference
of distinct points on a positive-genus curve
cannot be principal.

For any choice of rational $P$, the fixed module
$M_P^F$ therefore surjects onto a group of order
four and has nonzero classes represented by
non-point-theoretic full sections. All three
actual point sections are distinct by A1.
No all-level free-orbit property is claimed for
these nonzero module classes.


\paragraph{A fixed geometric curve defeats every finite prime list.}
Let $(C,P)$ be the explicitly specified characteristic-two
curve with its point at infinity in
Claim~\ref{clm:AGexplicitprimitivecounterexample}.
For every finite set $S$ of odd primes there is a finite
extension $L$ of its definition field such that no
$P$-split cyclic cover of degree $\ell\in S$ satisfies
(R) over $L$, or over any further finite extension.
Indeed, there are only finitely many geometric cyclic
covers of each fixed prime degree. Since no nontrivial
odd primitive cover of this curve has primitive two-rank
zero, the new abelian variety of each prime-degree cover
has a nonzero geometric two-torsion point. Choose one
such point for each cover. A common finite extension
makes all the covers and their deck groups defined,
makes a selected point above $P$ rational, and makes
all the selected new two-torsion points rational.

Over this field every geometric cover has its unique
$P$-split twist: constant cyclic twists are parametrized
by the Frobenius translation on the fibre over $P$, and
exactly the zero translation splits that fibre.
The chosen rational new two-torsion persists after every
further extension. Thus every cover in the finite list
has even relative class number throughout those
extensions. Conversely any $P$-split cover of one of
these degrees is the unique pointed form of a member
of the same finite geometric list, so none is omitted.
Consequently there is no finite list of witness primes
that works uniformly over all finite definition fields
even for this single geometric pointed curve. This
statement permits witnesses at other primes and does
not refute A3.
